\chapter{Results and Analysis}
\label{ch:results}

This chapter presents the implemented system and demonstrates, feature by feature,
how the requirements specified in Chapter~3 were realised in working software. Where
the preceding chapters reasoned about the system in the abstract---its architecture,
its data model and its planned behaviour---this chapter reports on the concrete
artefact that was built, deployed and exercised. The discussion follows the user's
journey through the deployed product---branded \textit{Sportify}---beginning with
onboarding and authentication, proceeding through the club, squad, competition and
match domains, and concluding with the medical, communication, analytics and
configuration modules. Each module is described along four consistent axes so that the
reader can move from \emph{intent} to \emph{mechanism} without losing the thread:
\begin{enumerate}
  \item the \emph{purpose} of the feature and the \emph{user experience} it delivers,
        illustrated by a screenshot of the running application;
  \item the \emph{back-end realisation}---which of the six Spring Boot
        microservices~\cite{springboot,microservices} owns the behaviour, the REST
        endpoints it exposes through the Spring Cloud Gateway~\cite{springcloud}, and
        the JPA entities and database tables that persist its state~\cite{postgres};
  \item the \emph{key logic and algorithms} in genuine detail---the fixture-generation
        routines, the standings recomputation, the live-match phase machine, the
        per-sport substitution rules, the in-match injury bridge, the status-driven
        medical workflow, the chat group/invitation model and WebSocket
        protocol~\cite{websocket}, the event-driven notification pipeline, the
        three-layer enforcement of role-based access~\cite{rbac}, and the Keycloak
        password re-verification~\cite{keycloak,jwt}; and
  \item the \emph{validation and edge-case handling} that makes each feature robust in
        the face of stale data, partial input and concurrent operation.
\end{enumerate}
The front-end is a Next.js~16 / React~19 application~\cite{nextjs,react} that talks to
the back-end exclusively through a single gateway origin, and every state-changing
action is persisted server-side so that a browser reload faithfully reconstructs the
application state. A substantial analytical discussion at the end of the chapter
(Section~\ref{sec:res-analysis}) evaluates how completely the delivered system
satisfies the functional and non-functional requirements, reports performance and
user-experience observations, and describes the testing approach that was followed.

%======================================================================
\section{Onboarding and Authentication}
\label{sec:res-login}

The entry point to the platform is a single sign-in surface that delegates identity
to Keycloak~\cite{keycloak}, the system's dedicated OpenID Connect identity provider.
Rather than implement its own credential store---an approach that would scatter
password hashing, lockout policy and token issuance across the domain services---the
platform treats identity as a cross-cutting concern owned by a purpose-built provider,
which is one of the architectural decisions defended in Chapter~3.

\paragraph{User experience.} A user submits a username and password on the sign-in
screen (Figure~\ref{fig:login}). On success the user is taken directly to the
operational home of the product; on failure a concise, non-revealing error is shown.
There is deliberately no separate landing page per role: every authenticated user
arrives at the single \texttt{/dashboard} route, whose contents are then computed from
the resolved role so that the same URL renders a different, permission-appropriate
surface for each user.

\paragraph{Back-end realisation.} The credentials are posted to the gateway, which
performs the OAuth~2.0 Resource Owner Password Credentials grant against the
\texttt{mscms} realm. Keycloak validates the credentials and returns a signed JSON Web
Token~\cite{jwt}. The token is the only proof of identity the rest of the system
trusts: every subsequent request carries it as a bearer token, the gateway validates
its signature and expiry before forwarding the request, and each microservice reads the
claims to authorise the operation. The user-management service (port~8081) brokers the
administrative side of identity---creating Keycloak users, assigning realm and client
roles, and reconciling the local user profile with the identity provider---through a
\texttt{KeycloakAdminService} that itself obtains a short-lived admin token from the
\texttt{master} realm before each privileged call.

\paragraph{Key logic.} The front-end decodes the \texttt{access\_token} client-side and
extracts two pieces of information. The token's \texttt{sub} claim is persisted as the
\texttt{keycloakId}---the canonical, stable identity used throughout the system, never
an internal database primary key, so that a service can correlate a player, a chat
member and an alert recipient without sharing a relational schema. The resolved role is
stored both in \texttt{localStorage}, where the API client reads it, and as a cookie
that the Next.js middleware reads on the server side to guard routes before a page is
ever rendered~\cite{rbac}. This dual storage is what makes the three-layer access
scheme described in Chapter~3 effective on the client tier: the cookie protects
navigation, while the token protects the API.

\paragraph{Validation and recovery.} Should a request later fail with HTTP~401---for
example after a token has expired, or after a database re-seed has invalidated the
previously issued token---the shared API client intercepts the response, clears local
storage and redirects the user transparently back to the login screen. This converts
what would otherwise be a confusing sequence of silently failing requests into a clean,
self-healing recovery path, and it guarantees that the application never operates in a
half-authenticated state.

\begin{figure}[htbp]
  \centering
  \includegraphics[width=0.85\textwidth]{fig-login.png}
  \caption{The Sportify sign-in screen. Credentials are validated against the Keycloak
  \texttt{mscms} realm through the API gateway, which returns a signed JWT that the
  front-end decodes to establish the user's identity and role.}
  \label{fig:login}
\end{figure}

%======================================================================
\section{Role-Aware Dashboard}
\label{sec:res-dashboard}

After authentication the user is presented with a unified dashboard that acts as the
operational home of the platform. The design intent is that a member of staff should be
able to assess the state of the club in a single glance, without first deciding which
of a dozen modules to open.

\paragraph{User experience.} Rather than maintaining a separate landing page per role,
the dashboard is assembled dynamically (Figure~\ref{fig:dashboard}). It composes a hero
banner carrying the club identity, a featured upcoming or most-recent match, a condensed
La~Liga standings strip, and a set of summary cards---squad size, scheduled matches,
active injuries and unread alerts---whose visibility is filtered by the permission
map~\cite{rbac}. A medical doctor and a sponsor consequently see overlapping but
distinct sets of cards, even though both navigated to the identical URL.

\paragraph{Back-end realisation.} The dashboard is a true aggregation surface: it has no
dedicated back-end of its own, but instead composes data fetched from several
microservices through the gateway on a single origin. Player counts are read from the
player-management service (port~8084), fixtures and live status from the training-match
service (port~8087), medical figures from the medical-fitness service (port~8086), and
unread-alert counts from the notification-mail service (port~8085). The fetch-based API
client issues these requests in parallel and the React layer composes the responses,
which keeps the perceived latency close to that of the slowest single call rather than
the sum of all of them.

\paragraph{Key logic and resilience.} Because the dashboard depends on several
independent services, it is designed to degrade gracefully: each summary card resolves
independently, so a momentarily unavailable medical service suppresses only the injuries
figure rather than blanking the whole page. The featured-match selection logic prefers
the next scheduled match and falls back to the most recent completed one, ensuring the
hero area is never empty even outside the playing season. This combination of parallel
aggregation and per-card isolation gives every member of staff an immediate,
role-appropriate and fault-tolerant snapshot of club operations on entry.

\begin{figure}[htbp]
  \centering
  \includegraphics[width=0.95\textwidth]{fig-dashboard.png}
  \caption{The role-aware dashboard. Summary cards, a featured match and a mini league
  table aggregate data from multiple microservices into a single overview surface, with
  card visibility governed by the user's role.}
  \label{fig:dashboard}
\end{figure}

%======================================================================
\section{Club Hub}
\label{sec:res-clubhub}

The Club Hub consolidates the modelled club's institutional identity into a single page.
Where the dashboard answers ``what is happening right now?'', the Club Hub answers ``who
are we?'', presenting the club as a continuous sporting institution rather than a set of
disconnected records.

\paragraph{User experience.} Taking FC~Barcelona as the working example, the Club Hub
presents the club crest and colours, the head coach, the home stadium, a trophy cabinet,
and the full squad rendered as player cards with real photographs
(Figure~\ref{fig:club-hub}). Crucially, the page also surfaces \emph{live} sporting
context rather than a static brochure: recent results are grouped by competition, a
``last-five'' form strip summarises momentum at a glance using a colour-coded
win/draw/loss sequence, and a \texttt{LIVE} badge appears whenever a match is in
progress.

\paragraph{Back-end realisation and the single-source-of-truth principle.} All of the
sporting data on this page originates from the back-end, which is the single source of
truth---the front-end never contacts the external football data provider directly. This
is a deliberate architectural constraint with several payoffs: the API key is never
exposed to the browser, the provider's rate limits are governed centrally, and the data
the user sees is exactly the data the rest of the system reasons about. Concretely, the
training-match service ingests fixtures and results on a schedule (see
Section~\ref{sec:res-laliga}) and the Club Hub reads the consolidated \texttt{Match}
records through the gateway, joining them client-side into the per-competition groupings
and the form strip.

\paragraph{Edge cases.} Because the modelled system date sits after the conclusion of
the imported seasons, the ``live'' badge and ``upcoming'' groupings are written to behave
correctly in the absence as well as the presence of in-progress matches: when no match is
live the badge is simply omitted, and when no fixtures remain the recent-results view is
shown in its place, so the page is always coherent regardless of where in the season the
club currently sits.

\begin{figure}[htbp]
  \centering
  \includegraphics[width=0.95\textwidth]{fig-club-hub.png}
  \caption{The Club Hub, presenting club identity, coaching and stadium details, the
  trophy cabinet, the squad, and results grouped by competition with a last-five form
  strip.}
  \label{fig:club-hub}
\end{figure}

%======================================================================
\section{Squad and Player Profiles}
\label{sec:res-players}

The Players module manages the multi-sport roster maintained by the player-management
service (port~8084). It is the module that most directly demonstrates the platform's
multi-discipline ambition.

\paragraph{User experience.} Players are listed as cards filterable by sport and
availability status, each carrying a real photograph, position and headline details
(Figure~\ref{fig:players}). Selecting a card opens an expanded detail view with the
player's biography, contract data and statistics, presented with a subtle ball-roll
animation that gives the transition a polished, product-quality feel
(Figure~\ref{fig:player-detail}).

\paragraph{Back-end realisation and data model.} Each player is a \texttt{Player} entity
keyed internally by a numeric identifier but correlated to identity by a
\texttt{keycloakId}, so the same person can be a squad member here, a chat participant in
the messaging service and an alert recipient in the notification service without any of
those services sharing a database. The entity carries a \texttt{photoUrl} populated from
public reference images, a \texttt{status} field (for example \texttt{AVAILABLE} or
\texttt{INJURED}), a \texttt{sportType} discriminator and the position and attribute
fields appropriate to the player's discipline. The module exposes the conventional REST
surface for listing, retrieving, creating and updating players, plus a dedicated
status-mutation endpoint that the live-match engine and the medical module call to flip a
player between availability states (Sections~\ref{sec:res-live}
and~\ref{sec:res-medical}).

\paragraph{Key logic.} Because the underlying model is sport-aware, the same listing
transparently accommodates footballers, basketballers, handball and tennis players, each
with the positions and attributes appropriate to their discipline, which directly
satisfies the multi-discipline requirement established in Chapter~2. The
\texttt{status} field is consumed across the platform as the authoritative statement of a
player's availability: a player marked \texttt{INJURED} is filtered out of lineup pickers,
match selection and training attendance automatically, which is what allows the medical
workflow to withdraw a player from operations simply by changing one field rather than by
editing every dependent module.

\begin{figure}[htbp]
  \centering
  \includegraphics[width=0.95\textwidth]{fig-players.png}
  \caption{The squad listing. Players are rendered as photograph-bearing cards,
  filterable by sport and availability, and served by the player-management service.}
  \label{fig:players}
\end{figure}

\begin{figure}[htbp]
  \centering
  \includegraphics[width=0.8\textwidth]{fig-player-detail.png}
  \caption{The expanded player profile, showing biography, contract and statistics for an
  individual squad member, opened with a subtle animation from the squad listing.}
  \label{fig:player-detail}
\end{figure}

%======================================================================
\section{Competitions}
\label{sec:res-competitions}

Competitions are modelled as first-class, database-backed entities owned by the
training-match service, in contrast to the throwaway, in-memory league tables of a
typical demonstration application. This decision is what allows the platform to enforce
realistic competition rules, to recompute standings deterministically, and to write live
results back to an authoritative fixture.

\paragraph{Data model.} Three persistent JPA entities capture the domain.
\texttt{Competition} records the competition's name, season, \texttt{type}
(\texttt{LEAGUE} or a knockout-style type) and the number of \texttt{rounds} a league
plays. \texttt{CompetitionTeam} records each participating team---its name, short name,
crest URL and sport---scoped to a competition. \texttt{CompetitionFixture} records a
single fixture: the home and away team identities and names, an optional \texttt{round}
label (empty for a plain league fixture, \texttt{LEAGUE\_STAGE} or a knockout label such
as \texttt{R16} otherwise), a \texttt{matchday} bucket, a \texttt{played} flag, the
recorded scores, and the timestamp at which the result was entered. These three entities
are exposed through a single controller mounted under \texttt{/matches/leagues} so that
no additional gateway route is required---the existing route to the training-match service
is reused, which keeps the gateway configuration small and auditable.

\paragraph{Scope and user experience.} The platform deliberately scopes football
competitions to the two that mirror the club's real calendar: \emph{La~Liga} (Spain's
Primera Divisi\'on, which the system treats as one and the same competition rather than
duplicating it) and the \emph{UEFA Champions League}. The Competitions landing page
presents these as selectable tabs and acts as the gateway into each competition's detailed
view (Figure~\ref{fig:competitions}). Real, completed seasons were imported as seed data
through a dedicated \texttt{/matches/leagues/import} endpoint that reads a season's matches
from the provider and materialises a competition, its teams and every finished fixture in
one transaction.

\paragraph{The read-only completed-competition rule.} A deliberate design decision marks
any competition whose every fixture has been played as \emph{read-only}: such a
competition displays a ``Completed---read-only'' badge and forbids result or date edits,
while an ongoing competition remains editable until its final result is recorded, at which
point it locks automatically. This rule exists to protect the integrity of historical
records---a finished season is a fact, not a draft---while still permitting an in-progress
competition to be updated match by match. Administrative deletion remains available in all
cases, since removing an entire competition is a coarse, intentional act rather than an
accidental edit.

\begin{figure}[htbp]
  \centering
  \includegraphics[width=0.95\textwidth]{fig-competitions.png}
  \caption{The Competitions landing page, with selectable tabs for La~Liga and the UEFA
  Champions League backed by persistent competition entities.}
  \label{fig:competitions}
\end{figure}

\subsection{La~Liga: Standings and Fixtures}
\label{sec:res-laliga}

The La~Liga view renders a live standings table together with the fixture and results
list for the season (Figure~\ref{fig:laliga}).

\paragraph{Standings computation.} The standings are not stored; they are recomputed from
the fixtures on every read, which guarantees they can never drift out of sync with the
underlying results. The computation iterates over the played fixtures, and for each one
credits three points to a win, one to each side of a draw, and zero to a loss, while
accumulating goals for, goals against and the played/won/drawn/lost tallies. The resulting
rows are then ordered by points, then by goal difference, then by goals scored---the
ordering used by the real competition---so that recording a single result and reloading the
table produces exactly the movement a spectator would expect. Because the recomputation is a
pure function of the fixture set, it is deterministic and trivially testable.

\paragraph{Ingestion pipeline.} For fixtures that are still to be played, the data is kept
current through the training-match service's ingestion pipeline. A
\texttt{FootballDataClient} together with a \texttt{MatchSyncService} polls the
football-data.org provider~\cite{footballdata} on a thirty-minute schedule---implemented as
a Spring \texttt{@Scheduled} task with an initial delay of roughly twenty-five seconds after
boot and a configurable interval---and on demand through an administrative
\texttt{/matches/sync-external} endpoint. The service reads the configured competitions
(\texttt{PD} for La~Liga and \texttt{CL} for the Champions League), filters the modelled
club's fixtures by the club's provider team identifier, and upserts them into the local
\texttt{Match} table. Each fixture is de-duplicated by an \texttt{externalId} derived from the
provider's match identifier, so repeated synchronisation is idempotent and never produces
duplicates.

\paragraph{Working within the provider's constraints.} The free tier of the provider forbids
team-scoped endpoints---a request to \texttt{/teams/\{id\}/matches} is rejected---so the
service instead reads the permitted competition-scoped match and standings endpoints and
filters client-side for the modelled club. The competition standings are served through a
read-through pass exposed at \texttt{/matches/standings/\{code\}} that is never persisted,
so the live league position always reflects the provider's current figures without the
platform attempting to mirror data it does not own. This design keeps the back-end
authoritative for the data the platform manages while still presenting accurate, real-world
league context, and it respects the provider's terms by centralising every external call
behind the API key held only on the server.

\begin{figure}[htbp]
  \centering
  \includegraphics[width=0.95\textwidth]{fig-laliga.png}
  \caption{The La~Liga competition view, showing the live standings table alongside the
  season's fixtures and results.}
  \label{fig:laliga}
\end{figure}

\subsection{Champions League: League Phase and Knockout Bracket}
\label{sec:res-cl-bracket}

The Champions League view reproduces the competition's modern format with exactly two tabs
(Figure~\ref{fig:cl-bracket}). The first presents the single \emph{league phase}, in which
all thirty-six participating teams accumulate points in one combined table; the fixtures of
this stage are stored under a \texttt{LEAGUE\_STAGE} round so that the engine can distinguish
them from knockout ties. The second presents the \emph{knockout bracket}, rendered as a
progression of ties from the knockout play-off and round of sixteen through to the final. The
knockout tab remains locked until the league phase is complete, mirroring the real
competition's sequencing and preventing the user from advancing prematurely.

\paragraph{Why the round label matters.} The same \texttt{CompetitionFixture} model underpins
both stages, with each fixture's \texttt{round} value classifying it as either a points fixture
or a knockout tie. This single discriminator carries real behavioural weight: as
Section~\ref{sec:res-live} describes, the live-match engine consumes it to decide whether a
draw is a valid result. A league-phase fixture permits a draw and the engine stops at full
time; a knockout tie does not, and the engine extends into extra time and, if necessary, a
penalty shootout. Threading the format rule through one authoritative field, rather than
duplicating it in the user interface, is what keeps the live engine and the competition record
consistent.

\begin{figure}[htbp]
  \centering
  \includegraphics[width=0.95\textwidth]{fig-cl-bracket.png}
  \caption{The UEFA Champions League knockout bracket, presenting the progression of ties
  from the round of sixteen to the final. The bracket is unlocked only once the league
  phase has concluded.}
  \label{fig:cl-bracket}
\end{figure}

\subsection{Creating a Competition}
\label{sec:res-create-competition}

Administrators may author a new competition through a guided, multi-step flow
(Figure~\ref{fig:create-competition}). The user first selects a preset (La~Liga or
Champions League), then a season from a curated list, and finally the participating teams:
La~Liga is fixed at its twenty members, whereas the Champions League flow lets the
administrator select thirty-six clubs from a curated European team pool.

\paragraph{La~Liga generation: the double round-robin.} For a league competition the
back-end pre-generates the entire schedule on creation so that the user can only ever enter
results for genuine fixtures and can never inflate the table by recording the same pairing
repeatedly. The generator implements a double round-robin: it iterates over the configured
number of rounds (two for La~Liga) and, within each round, over every unordered pair of
teams, emitting one fixture per pair. The home and away sides are swapped on alternate rounds
so that each pairing is played once at home and once away, exactly as the real competition is
structured, and each fixture is assigned an incrementing \texttt{matchday} for display
grouping. For twenty teams this yields the expected $20 \times 19 = 380$ fixtures.

\paragraph{Champions League generation: the league-phase rotation algorithm.} The Champions
League league phase is the most algorithmically interesting case, because each of the
thirty-six teams must play exactly eight \emph{distinct} opponents---it is not a full
round-robin, in which every team would play every other. The platform generates this with a
deterministic circle/rotation pairing. The algorithm maintains, for every team, a running
\emph{degree} (the number of opponents already assigned) and a set of pair keys already used.
It then sweeps over rotational offsets: for an offset $o$ it considers pairing team $i$ with
team $(i + o) \bmod n$, and accepts the pairing only if both teams still have spare degree
and the unordered pair has not already been scheduled. Home and away are assigned by a parity
rule so that home counts stay balanced across the field. The sweep continues, increasing the
offset, until the total number of fixtures reaches $\tfrac{n \times 8}{2} = 144$, at which
point every team has exactly eight distinct opponents and no pairing is repeated. The
fixtures are then distributed across eight matchday buckets purely for display grouping, and
every fixture is tagged with the \texttt{LEAGUE\_STAGE} round so that it renders in the league
phase and is correctly treated as draw-permitting by the live engine. Because the algorithm is
deterministic, the same team list always produces the same regular, well-spread fixture graph,
which makes the generated schedule reproducible and verifiable.

\paragraph{Persistence and atomicity.} On confirmation the back-end generates the appropriate
fixture set---the round-robin for a league, or a bulk insert of the rotation-generated
league-phase fixtures for the Champions League via a \texttt{/matches/leagues/\{id\}/fixtures/bulk}
endpoint---and persists the competition together with its teams and fixtures inside a single
transaction. This converts what would traditionally be a manual, error-prone, spreadsheet-driven
exercise into one auditable, all-or-nothing operation: either the entire competition is created
correctly or none of it is.

\begin{figure}[htbp]
  \centering
  \includegraphics[width=0.9\textwidth]{fig-create-competition.png}
  \caption{The guided competition-creation flow: preset selection, season selection and
  team selection culminate in a generated fixture set persisted by the training-match
  service.}
  \label{fig:create-competition}
\end{figure}

%======================================================================
\section{Head-to-Head Records}
\label{sec:res-h2h}

The Head-to-Head view answers a question that fans and analysts ask constantly but that is
tedious to compute by hand: how has the club fared against a given opponent over time?

\paragraph{User experience.} The page lists every opponent the club has met, and for each one
summarises the all-time record---matches played, won, drawn and lost, together with goals
scored and conceded (Figure~\ref{fig:h2h}). Selecting an opponent expands the individual
fixtures that make up the record, so a headline statistic can be traced back to the matches
that produced it.

\paragraph{Back-end realisation and logic.} The view is computed entirely from the
authoritative \texttt{Match} records held by the training-match service, the same records that
the ingestion pipeline populates and that the live engine writes back to. The front-end groups
the club's matches by opponent and folds each group into a win/draw/loss and goals-for/against
summary using the same scoring conventions as the standings computation, so a head-to-head row
is always consistent with the competition tables. Because it derives from stored results rather
than a separate data source, the record updates automatically as new matches complete---no
manual maintenance of a head-to-head ledger is required.

\begin{figure}[htbp]
  \centering
  \includegraphics[width=0.95\textwidth]{fig-head-to-head.png}
  \caption{The Head-to-Head view, summarising the club's all-time record against each opponent
  (played, won, drawn, lost and goals) computed from the authoritative match records, with the
  underlying fixtures expandable for any opponent.}
  \label{fig:h2h}
\end{figure}

%======================================================================
\section{World Cup}
\label{sec:res-worldcup}

The World Cup view brings a major international tournament into the platform as live
sporting context, demonstrating that the competition machinery generalises beyond the
club's own two domestic competitions.

\paragraph{User experience.} The page presents the tournament in its real structure: a set
of group tables followed by a knockout bracket that progresses to the final
(Figure~\ref{fig:worldcup}). Group standings and knockout fixtures are shown with the same
visual language as the club's own competitions, so a user moves between the domestic and
international views without re-learning the interface.

\paragraph{Back-end realisation.} As with every external dataset, the front-end never calls
the provider directly. The training-match service exposes read-through endpoints for the
tournament---one that returns the group tables and one that returns the tournament's
matches---so the World Cup competition is available on the provider's free tier and is served
to the browser through the gateway. The group tables are computed and the bracket assembled on
the server side, keeping the browser a pure presentation layer and the API key confined to the
back-end.

\paragraph{Edge cases.} The view is written to render correctly at every phase of the
tournament: before the knockout stage begins it shows only the group tables, and as ties are
decided the bracket fills in progressively, so the page is coherent whether the tournament is
in its group phase, its knockout phase, or complete.

\begin{figure}[htbp]
  \centering
  \includegraphics[width=0.95\textwidth]{fig-world-cup.png}
  \caption{The World Cup view, presenting the tournament's group tables and knockout bracket.
  Group standings and the bracket are computed on the server and served to the browser through
  the gateway via read-through endpoints, so the external data provider is never called from
  the client.}
  \label{fig:worldcup}
\end{figure}

%======================================================================
\section{Trophy Room}
\label{sec:res-trophy}

The Trophy Room is the celebratory counterpart to the operational modules: it presents the
club's honours as a curated cabinet, reinforcing the institutional identity that the Club Hub
introduces.

\paragraph{User experience.} The page displays the club's trophies grouped by competition and
era, each rendered as a distinct honour with its competition and the seasons in which it was
won (Figure~\ref{fig:trophy}). The layout is intentionally visual and museum-like, contrasting
with the dense tabular surfaces elsewhere in the product, and gives the platform a sense of
heritage rather than presenting the club as a purely transactional set of records.

\paragraph{Realisation.} The honours are presented as part of the club-identity surface,
complementing the trophy summary shown on the Club Hub with a dedicated, expanded gallery. By
giving the achievements their own room, the platform acknowledges that a club is defined as much
by its history as by its current fixtures, which strengthens the product's claim to be a
faithful digital home for the institution rather than a back-office tool.

\begin{figure}[htbp]
  \centering
  \includegraphics[width=0.9\textwidth]{fig-trophy-room.png}
  \caption{The Trophy Room, presenting the club's honours as a curated cabinet grouped by
  competition and era, complementing the trophy summary on the Club Hub.}
  \label{fig:trophy}
\end{figure}

%======================================================================
\section{Match Hub}
\label{sec:res-matches}

The Match Hub is the operational list of the club's own matches---the matches it schedules,
plays and reviews, as distinct from the external fixtures surfaced in the competition views.

\paragraph{User experience.} It comprises a \emph{Schedule a match} action, four status filters
(Scheduled, Live, Completed and a catch-all) and a team filter, with each match rendered as a
card showing the opponent's name and crest, the kickoff time, and---for upcoming matches---a
live countdown that flips the card to a \texttt{LIVE} state automatically when kickoff arrives
(Figure~\ref{fig:matches}). The countdown is computed entirely client-side from the stored
kickoff timestamp, so a card transitions to live without any server round-trip and without the
user refreshing the page.

\paragraph{Self-cleaning behaviour.} The Hub also performs housekeeping on load to keep the list
trustworthy without manual intervention. Matches that were scheduled or marked live but have
passed kickoff by more than ninety minutes with no recorded events are treated as abandoned and
removed---these are typically test or mistakenly created matches that were never played. Matches
that carry real events but were never formally ended are finalised to a completed state. To make
finished matches self-consistent, goal events are backfilled so that each completed match's
recorded scorers reconcile with its final score, drawing the scorers from the saved starting XI.
The net effect is that the Match Hub always presents a clean, internally consistent record, and a
stray or interrupted match cannot accumulate as noise.

\begin{figure}[htbp]
  \centering
  \includegraphics[width=0.95\textwidth]{fig-matches.png}
  \caption{The Match Hub, listing the club's matches with status and team filters and a
  schedule action. Upcoming cards display a live countdown and transition to a LIVE state
  at kickoff.}
  \label{fig:matches}
\end{figure}

\subsection{Scheduling a Match}
\label{sec:res-match-schedule}

Scheduling distinguishes between a \emph{Friendly} and an \emph{Official} match
(Figure~\ref{fig:match-schedule}). A Friendly allows a free choice of opponent, whereas an
Official match is bound to a real, still-running competition and offers only that competition's
unplayed fixtures, carrying the fixture's identity forward so that the eventual result can be
written back to the source. This binding is what closes the loop between the live engine and the
competition record: an official match is not an island, but an instance of a specific
competition fixture.

\paragraph{Kickoff validation.} The kickoff time is validated to lie at least five minutes in the
future. The validation is a pure function of the chosen timestamp and the current time: it rejects
an empty value, an unparseable value, a value in the past, and a value less than five minutes
ahead, returning a specific user-facing message for each case. The picker itself advertises a
minimum selectable value of ``now plus five minutes'', nudging the user toward a valid choice
before they even submit. A subtle but important detail is that the kickoff is stored as a naive
local wall-clock string rather than being passed through a timezone conversion: an earlier defect
in which times were shifted by the local UTC offset (making matches appear several hours early and
spuriously live) was fixed by normalising the picker value to the back-end's naive
\texttt{LocalDateTime} format without any UTC arithmetic.

\paragraph{Lineup is mandatory.} Critically, a match is never persisted on its own: a lineup is
required. The scheduling form stages its payload and hands off to the lineup builder, and only when
a formation and starting eleven are confirmed does the system create the match, its formation and
its lineup together in one coherent operation. This guarantees a key invariant exploited everywhere
downstream---every match in the system, the moment it exists, already has a reconstructable starting
XI.

\begin{figure}[htbp]
  \centering
  \includegraphics[width=0.9\textwidth]{fig-match-schedule.png}
  \caption{Scheduling a match. The Friendly/Official chooser, opponent or fixture
  selection and a future-dated kickoff feed a draft that is only persisted once a lineup
  has been confirmed.}
  \label{fig:match-schedule}
\end{figure}

\subsection{The Lineup Builder}
\label{sec:res-lineup}

The lineup builder is a formation-first, interactive pitch (Figure~\ref{fig:lineup}).

\paragraph{User experience.} The user begins by selecting a formation appropriate to the sport,
after which players are placed and rearranged on the pitch by drag-and-drop, with each position
rendered as a card bearing the player's photograph. A single-card drag ghost gives the drag
operation a clean, unambiguous visual, and players can be swapped between positions or between the
pitch and the bench. The starting eleven and the bench are assembled here.

\paragraph{Back-end realisation and data model.} The builder is backed by two entities in the
training-match service. A \texttt{MatchFormation} captures the chosen shape and a compact
\texttt{formationDetails} description, while \texttt{MatchLineup} records the per-player placements.
These are exposed through \texttt{/match-formations} and \texttt{/match-lineups} endpoints and are
saved atomically with the match. A practical constraint worth recording is that the
\texttt{formationDetails} column is a bounded \texttt{VARCHAR(255)}, so the encoding is kept compact
to avoid an over-length write being rejected---a small but real example of the data model shaping
the client encoding.

\paragraph{Why this matters.} Because the lineup is captured at creation time and locked thereafter,
every completed match in the system carries a genuine, reconstructable starting XI rather than a
placeholder. This is precisely what makes the in-match event and substitution logic meaningful: a
substitution refers to real players who were really on the pitch, an in-match injury flags a real
squad member, and the match-details page can rebuild the exact eleven that started. The entire fifty-odd
finished football matches in the seed data were backfilled with real starting elevens for the same
reason, so that the historical record is as self-consistent as a freshly scheduled match.

\begin{figure}[htbp]
  \centering
  \includegraphics[width=0.95\textwidth]{fig-lineup.png}
  \caption{The formation-first lineup builder. Players are positioned on the pitch by
  drag-and-drop, and the resulting formation and starting XI are persisted atomically with
  the match.}
  \label{fig:lineup}
\end{figure}

%======================================================================
\section{The Live-Match Engine}
\label{sec:res-live}

The live-match engine is the most behaviourally rich component of the platform, and the one in
which the front-end and back-end cooperate most tightly. It transforms a static fixture into a
real-time, broadcast-style operational surface (Figure~\ref{fig:live-match}).

\subsection{Gating and Auto Go-Live}
Control of the engine is gated until kickoff is reached or the match is already live, so a match
cannot be operated before it is due to begin. The transition itself is automatic: pure timestamp
helpers compute, from the stored kickoff and the current time, whether kickoff has arrived, the
running match minute, and the phase the minute implies. When the kickoff moment passes, the match
goes live on its own---the user never manually flips it---and the elapsed minute begins advancing.
Because these computations are pure functions of two timestamps, the common case of simply watching
the clock advance requires no server round-trip, which keeps the live view responsive.

\subsection{The Phase Machine}
The engine drives a broadcast-style \emph{phase machine}. Each phase is described by a small record
giving its label, the base minute at which its clock starts, its regulation length, whether it is a
playing phase with a running clock, and the phase that follows it. The regulation sequence is
\texttt{FIRST\_HALF}~(45') $\rightarrow$ \texttt{HALF\_TIME} $\rightarrow$
\texttt{SECOND\_HALF}~(45') $\rightarrow$ \texttt{FULL\_TIME}. For a true knockout tie that remains
level at full time, the machine extends into \texttt{EXTRA\_FIRST}~(15') $\rightarrow$
\texttt{EXTRA\_HALFTIME} $\rightarrow$ \texttt{EXTRA\_SECOND}~(15') and, if still level, a
\texttt{PENALTIES} shootout, which is a non-timed state owned explicitly by the user interface. For a
points fixture---including a Champions League league-phase match---a draw is a legitimate result and
the engine stops at full time.

\paragraph{Resolving knockout from points authoritatively.} The knockout-versus-points decision is
the linchpin of the whole machine, and it is resolved authoritatively rather than guessed. A set of
match types (league, friendly, group and the various league-stage labels) are treated as
points fixtures for which a draw is valid; but when a match originates from a real competition fixture,
the resolved round of that fixture is checked \emph{first} and overrides the loose type. The effect is
that a match stored loosely as a tournament match is still correctly treated as a draw-permitting league
fixture when its round indicates the league phase, and a genuine knockout tie correctly extends into
extra time and penalties. Threading the decision through the originating fixture's round is what keeps
the engine consistent with the competition format described in
Section~\ref{sec:res-cl-bracket}.

\subsection{The Clock, Stoppage Time and Phase-Bounded Minutes}
Each playing phase exposes an administrator-set stoppage time, allowing the clock to read ``45+3''' or
``90+5''' during added time. The clock is carefully bounded so that it can never run away into nonsense.
The displayed minute is hard-capped at the phase's regulation length, and the added portion shown is the
\emph{admin-set} stoppage only, never the raw elapsed overflow. Consequently a half left running for an
extended period freezes at, for example, ``45+3''' (or simply ``45''' when no stoppage was set) rather
than inflating into an absurd minute---a concrete defect (an observed ``45+132''') that was fixed by this
capping rule. It is always the user who advances the phase; the clock follows, it does not lead.

Event entry is likewise phase-bounded. The recordable-minute range is computed per phase: the first half
accepts minutes one to forty-five plus any stoppage, the second half forty-six to ninety plus stoppage,
and each extra-time period its fixed fifteen-minute window. Non-playing phases (the half-time and full-time
intervals, the extra-time interval and the penalty shootout) carry no recordable minute. This prevents a
nonsensical entry such as a first-half goal recorded at minute eighty, and it means that when state is
reconstructed on reload, a stored event minute maps unambiguously back to the phase in which it occurred.

\subsection{Events, Scoring and Substitutions}
Events are entered against the live clock and persist immediately through the match-events API. A
\texttt{MatchEvent} entity records the event type, the minute, the acting player's \texttt{keycloakId} and
a free-text description; the event-type enumeration is broad and sport-aware, covering goals, assists,
penalties scored and missed, own goals, yellow and red cards, corners, free kicks, offsides and VAR reviews
for football, as well as the equivalent events for basketball, handball, tennis, volleyball and swimming.
Scoring events increment the appropriate side's tally, and the running score is written back to the match as
play proceeds.

Substitutions have a dedicated control that records the outgoing and incoming players and the minute. Because
the underlying \texttt{MatchEvent} carries a single player reference, a substitution encodes both players---their
numeric identifiers, their \texttt{keycloakId}s and their names---in a parseable tag inside the event description,
which lets the timeline render both faces and lets the post-match review form learn which substitutes actually
came on. Substitutions are limited per sport: football permits the modern maximum of five, while free-interchange
roster sports such as basketball and handball are effectively unlimited. The engine tracks which players have left
the pitch, which substitutes have entered, and how many of the allowance has been consumed, and it rebuilds this
on-pitch set on reload from the stored events so that the live state survives a page refresh.

\subsection{The In-Match Injury Bridge}
An in-match \texttt{INJURY} event does double duty and is the clearest demonstration of genuine cross-service
integration. It is recorded as an ordinary match event \emph{and} simultaneously creates a real injury record in
the medical-fitness service: the engine posts a new injury (in the \texttt{REPORTED} status, linked to the player,
with a chosen injury type and severity) and then flips the affected player's status to \texttt{INJURED} through the
player-management API. The consequences ripple outward without any further action: the injured player is forced off
the pitch and, because an injury substitution is exempt from the normal sub-count, the replacement is always
permitted even past the five-substitution limit; the player becomes unselectable in future lineups, match selection
and training; and the incident appears in the Medical module's Treatment Room ready for the staged clinical journey
described in Section~\ref{sec:res-medical}. A single click in the live engine thus drives three services at once.

\subsection{Alerts and Result Write-Back}
Recording a goal raises a real, persisted alert so that interested roles are notified as the match unfolds, and
ending the match emits a result-available alert and writes the final score back to the source competition fixture.
The alert posting is best-effort and resilient: the engine attempts a goal-scored alert type and, because the
back-end alert enumeration must be respected, falls back to a valid match-scoped type if necessary so that a real
alert is always created and the call never throws. Ending an official match writes the final score to its originating
fixture, which in turn moves the competition standings (Section~\ref{sec:res-laliga}) and, where relevant, unlocks or
advances the knockout bracket. This closes the loop between live operations, the notification fabric and the
competition record, so that operating a live match is not a self-contained activity but one that updates the whole
platform.

\begin{figure}[htbp]
  \centering
  \IfFileExists{screenshots/fig-live-match.png}{\includegraphics[width=0.95\textwidth]{fig-live-match.png}}{\fbox{\parbox[c][5cm][c]{0.9\textwidth}{\centering\itshape Live-match control panel --- screenshot to be inserted (\texttt{fig-live-match.png}).}}}
  \caption{The live-match engine. The broadcast-style phase machine, stoppage-time
  control, event entry, per-sport substitutions and in-match injury recording drive the
  match in real time and persist every action to the back-end.}
  \label{fig:live-match}
\end{figure}

\subsection{Match Details: Events, Statistics and Reviews}
\label{sec:res-match-details}

Once a match is complete its detail page becomes a read-only record of what occurred
(Figure~\ref{fig:match-details}). It reconstructs the saved starting XI on a pitch board, presents a
chronological events timeline, a statistics panel, and a reviews wall on which coaching staff record
post-match player assessments. The reviews wall reads the substitution events to learn which substitutes
came on, so an assessment can be recorded for a player who only featured for the closing minutes, not merely
for the eleven who started. Because every finished match carries a real lineup and a backfilled set of goal
events consistent with its score, the detail page presents a faithful, self-consistent account of the match
rather than a synthetic summary, and a finished match is rendered read-only in keeping with the platform-wide
treatment of completed records.

\begin{figure}[htbp]
  \centering
  \includegraphics[width=0.95\textwidth]{fig-match-details.png}
  \caption{The match details page, presenting the reconstructed starting XI, the events
  timeline, match statistics and the post-match player-review wall for a completed match.}
  \label{fig:match-details}
\end{figure}

%======================================================================
\section{Medical and Fitness}
\label{sec:res-medical}

The medical module, served by the medical-fitness service (port~8086), manages the full clinical lifecycle
of a player's injury. It is one of the modules that most clearly transcends simple record-keeping, encoding a
genuine clinical workflow rather than a flat table of incidents.

\paragraph{The Treatment Room.} The module's centrepiece is the \emph{Treatment Room}, which divides the squad
into \emph{Currently Injured} and \emph{Recovered} cohorts (Figure~\ref{fig:medical-room}). A player is
considered injured---and is automatically withdrawn from upcoming matches, lineups and training---from the moment
an injury is logged until a passing fitness test clears them back to availability. This withdrawal is achieved
not by editing every dependent module but by the player's single \texttt{status} field, which every other module
already respects, so the squad's operational status stays consistent across the match and training modules without
any manual reconciliation.

\paragraph{Back-end realisation and data model.} An \texttt{Injury} entity records the affected player, the body
part, the injury type and severity, the report metadata, and---most importantly---a \texttt{status} drawn from a
small enumeration: \texttt{REPORTED}, \texttt{DIAGNOSED}, \texttt{TREATING}, \texttt{RECOVERING},
\texttt{RECOVERED} and \texttt{CHRONIC}. The injury aggregates its associated clinical sub-records: diagnoses,
treatments and rehabilitation entries, with recovery programmes and fitness tests recorded as the journey
advances. The injury exposes a conventional REST surface for creation and retrieval, and it is this entity that the
live-match engine targets when an in-match injury occurs.

\begin{figure}[htbp]
  \centering
  \includegraphics[width=0.95\textwidth]{fig-medical-room.png}
  \caption{The Medical Treatment Room, separating currently injured players from those who
  have recovered. Logging an injury automatically withdraws the player from matches,
  lineups and training until a passing fitness test restores availability.}
  \label{fig:medical-room}
\end{figure}

\subsection{The Injury Journey}
\label{sec:res-medical-journey}

Each injured player has a staged recovery journey rendered as a stepper:
Injury~$\rightarrow$~Diagnosis~$\rightarrow$~Treatment~$\rightarrow$~Rehabilitation~$\rightarrow$~Recovery~$\rightarrow$~Fitness~Test~$\rightarrow$~Recovered
(Figure~\ref{fig:medical-journey}).

\paragraph{Status as the single source of truth.} The injury's own status is the single source of truth for how far
the player has progressed, and this is the design decision that makes the workflow robust. Rather than inferring the
current stage from whichever sub-records happen to exist---an approach that breaks down the moment a stray document
is left over from a previous episode---the stepper reads the injury's status directly and maps it onto the stage. As
the player progresses, the status advances through the enumeration: an initial report is \texttt{REPORTED}, a recorded
diagnosis moves it to \texttt{DIAGNOSED}, treatment to \texttt{TREATING}, and so on through \texttt{RECOVERING} to
\texttt{RECOVERED}. Because the stage is read from one authoritative field, the stepper always reflects the true
clinical state and is immune to leftover sub-records, which was a concrete source of confusion before the status-driven
design was adopted.

\paragraph{Ordered progression.} Advancing a stage records a new clinical document---a diagnosis, a treatment, a
rehabilitation plan, a recovery programme or a fitness-test result---and the journey enforces the correct ordering. It
guarantees, for example, that a rehabilitation record exists before a recovery programme can be created, so the clinical
history reads as a coherent sequence rather than an arbitrary set of documents. Only a \emph{passing} fitness test
completes the journey: it moves the injury to recovered and returns the player to the available pool, at which point the
player reappears in lineup pickers, match selection and training. A failing fitness test does not clear the player,
keeping a not-yet-fit player correctly withheld from selection.

\begin{figure}[htbp]
  \centering
  \includegraphics[width=0.95\textwidth]{fig-medical-journey.png}
  \caption{The per-player injury journey stepper, progressing from the initial injury
  through diagnosis, treatment, rehabilitation, recovery and a fitness test to full
  recovery. The injury's status is the single source of truth for the current stage.}
  \label{fig:medical-journey}
\end{figure}

\subsection{General Medical Catalogs}
\label{sec:res-medical-catalog}

Alongside the per-player journey, the module exposes general catalogs for each clinical category---Injuries, Diagnoses,
Treatments, Rehabilitation, Recovery and Fitness Tests (Figure~\ref{fig:medical-catalog}). These provide reusable,
browsable records and serve as templates that the journey's stage pickers can draw upon.

\paragraph{Copy, not link.} A subtle but deliberate design choice governs how a catalog entry is used: selecting one
pre-fills the relevant stage form, but its content is then \emph{copied} into a fresh record for the current player rather
than re-linked to the shared template. This preserves a clean, independent per-injury history---editing a player's
treatment record cannot retroactively mutate the shared catalog, and the catalog cannot change a player's recorded
history---while still saving the clinician the effort of re-entering common information. The result combines the
convenience of templates with the integrity of immutable per-episode records.

\begin{figure}[htbp]
  \centering
  \includegraphics[width=0.95\textwidth]{fig-medical-catalog.png}
  \caption{The general medical catalogs. Reusable injury, diagnosis, treatment,
  rehabilitation, recovery and fitness-test records double as templates that pre-fill the
  journey's stage forms.}
  \label{fig:medical-catalog}
\end{figure}

%======================================================================
\section{Team Communication}
\label{sec:res-messages}

The Messages module delivers a real-time, group-based team chat backed by the notification-mail service
(port~8085). It gives the platform a self-contained communication channel without depending on any third-party
messaging service.

\paragraph{Group and invitation model.} The data model comprises a \texttt{ChatGroup} and a
\texttt{ChatGroupMember}, the latter carrying a membership status (\texttt{INVITED} or \texttt{MEMBER}) that
distinguishes an invited user from a confirmed member. A seeded general group is returned to everyone, while a
user-created group makes its creator a member and places invitees in the invited state until they accept. The chat
group controller is mounted at \texttt{/messages/groups}---a literal path chosen so it takes precedence over the
parameterised \texttt{/messages/\{id\}} route---and exposes the full lifecycle: listing groups, fetching a group and
its members, creating a group, listing a user's pending invitations, and accepting, declining or issuing invitations
through dedicated \texttt{/accept}, \texttt{/decline} and \texttt{/invite} endpoints. The \texttt{Message} entity
carries a \texttt{groupId}, scoping every message to its conversation, and supports threaded replies through a parent
reference.

\paragraph{Reactions and de-duplication.} The conversation supports emoji reactions and client-side de-duplication so
that each message renders exactly once even when it arrives by more than one path (for instance, both echoed over the
socket and returned by a poll). De-duplication is layered: a message is keyed first by its server identifier, then by a
client-assigned message identifier, and finally by the combination of sender and content, which together resolved an
earlier double-render defect.

\paragraph{Real-time delivery over WebSocket.} Delivery is genuinely real-time. A native browser WebSocket connects
directly to the notification service at \texttt{ws://localhost:8085/ws/chat}~\cite{websocket}. The path sits outside the
\texttt{/messages} namespace and outside the gateway deliberately: the raw socket bypasses the gateway's JWT filter and
connects straight to the service, where the endpoint is permitted, which avoids the complications of authenticating a
long-lived socket through the HTTP gateway. The protocol is a small set of JSON text frames. A client first sends an
\texttt{identify} frame carrying its identity; thereafter it sends \texttt{chat} frames carrying the group, sender and
content, which the server re-broadcasts verbatim to every connected client; a \texttt{ping} frame keeps the connection
alive; and the server emits a \texttt{presence} frame on every connect, identify and disconnect, reporting how many users
are online. HTTP polling is retained as a fallback should the socket be unavailable, so the chat continues to function---if
less instantly---even when a WebSocket connection cannot be established (Figure~\ref{fig:messages}).

\begin{figure}[htbp]
  \centering
  \includegraphics[width=0.95\textwidth]{fig-messages.png}
  \caption{The team group chat. Group conversations, invitations, threaded replies and
  emoji reactions are delivered over a native WebSocket to the notification-mail service,
  with polling as a fallback.}
  \label{fig:messages}
\end{figure}

%======================================================================
\section{Notifications and Alerts}
\label{sec:res-notifications}

Notifications are event-driven and surfaced through a bell in the top navigation bar (Figure~\ref{fig:notifications}).
The notification fabric is one of the clearest expressions of the platform's event-driven, loosely coupled design.

\paragraph{The event pipeline.} The notification-mail service consumes domain events from Kafka~\cite{kafka} through a
set of \texttt{@KafkaListener} consumers, one per topic. When the medical service reports an injury, the player service
completes a transfer, the training service cancels a session, or the match service schedules or completes a match, the
producing service emits an event onto the corresponding Kafka topic following the outbox pattern, and the notification
service's \texttt{EventConsumer} materialises it as a targeted alert with an appropriate type and priority. An injury
event, for instance, becomes an \texttt{INJURY\_REPORTED} alert routed to medical staff, with the priority derived from
the injury's severity; a transfer event becomes a \texttt{TRANSFER\_COMPLETED} alert and may additionally trigger an email
to the affected player. The live-match engine contributes alerts directly over HTTP rather than through Kafka, posting a
goal alert as the match unfolds and a result-available alert when it ends. The alert type enumeration is correspondingly
broad, spanning medical, performance, match-and-training, player-management and scouting-and-sponsorship categories.

\paragraph{The notification centre.} The bell opens a notification centre that presents alerts across two tabs,
distinguishes alerts from notifications, carries unread badges, and refreshes on a short polling interval (on the order of
a few seconds) so that critical events reach the relevant actor promptly. Competition fixtures deep-link from the
notification and competition surfaces into the match-creation flow, pre-filling the teams and carrying the originating
fixture's identity so that scheduling a match from an alert requires no re-entry of detail. This realises the role-aware
notification fabric called for in the requirements and removes the reliance on ad-hoc, manual alerting.

\paragraph{Why event-driven.} Routing notifications through Kafka rather than direct service-to-service calls is what keeps
the producing services ignorant of who consumes their events. The medical service does not know that an injury produces an
alert, an email and a player withdrawal; it simply records the injury and emits a fact. New consumers can be added without
touching the producers, which is the loose coupling the microservice architecture set out to achieve.

\begin{figure}[htbp]
  \centering
  \includegraphics[width=0.9\textwidth]{fig-notifications.png}
  \caption{The notification centre opened from the top-bar bell. Alerts are generated
  automatically from Kafka domain events and from the live-match engine, tagged by role and
  unread state.}
  \label{fig:notifications}
\end{figure}

%======================================================================
\section{Training}
\label{sec:res-training}

The Training module presents the club's training schedule as managed by the training-match service, covering sessions,
plans, drills, attendance and post-session assessments (Figure~\ref{fig:training}).

\paragraph{User experience.} Sessions are displayed on a schedule with a \emph{live-now} indicator that highlights a
session currently in progress, computed from the session's scheduled window in the same timestamp-driven manner as the
match countdown. Plan-driven training ties drills to ratings, attendance and completion tracking, so a session is not a
bare calendar entry but a structured plan whose execution can be recorded: which drills were run, who attended, how each
drill was rated and whether the plan was completed.

\paragraph{Integration with availability.} Because the medical module withdraws injured players from training
automatically---again through the shared player status field---the training schedule always reflects the squad's true
availability, and an injured player cannot be marked present at a session he is medically barred from. When a session is
cancelled, the event flows onto Kafka and the notification fabric raises a training-cancelled alert, so the cancellation
reaches the affected staff and players without a separate announcement step. This connects training to the medical and
notification domains rather than leaving it as an isolated calendar.

\begin{figure}[htbp]
  \centering
  \includegraphics[width=0.95\textwidth]{fig-training.png}
  \caption{The training schedule with a live-now indicator, showing plan-driven sessions
  with drills, attendance and completion tracking.}
  \label{fig:training}
\end{figure}

%======================================================================
\section{Scouting}
\label{sec:res-scouting}

The Scouting module supports the club's recruitment workflow, allowing scouts to record prospect reports against
detailed scouting fields and to route findings to the sport manager for consideration (Figure~\ref{fig:scouting}).

\paragraph{Realisation and integration.} A scouting report captures structured attributes of a prospect rather than
free prose, which makes reports comparable and searchable instead of opaque. When a report is submitted it can raise a
scout-report alert to the sport manager through the same notification fabric, so a finding does not languish unread in an
inbox. Reports feed the transfer and analytics surfaces, providing the audit trail that traditional email-and-PDF scouting
processes lack and connecting the recruitment pipeline to the rest of the platform: a prospect assessed by a scout, transferred
by the player service and analysed by the reports service is the same identity throughout, correlated by the canonical
\texttt{keycloakId}.

\begin{figure}[htbp]
  \centering
  \includegraphics[width=0.95\textwidth]{fig-scouting.png}
  \caption{The Scouting module, where scouts record prospect reports against structured
  fields for review by the sport manager.}
  \label{fig:scouting}
\end{figure}

%======================================================================
\section{Sponsors}
\label{sec:res-sponsors}

The Sponsors module digitises commercial negotiations (Figure~\ref{fig:sponsors}).

\paragraph{Realisation.} Sponsor offers and contract terms are captured as structured records that move through a
negotiation workflow, replacing the opaque email threads of traditional practice with an auditable, notification-generating
process. An incoming offer can raise a sponsor-offer alert to the relevant commercial roles, and decisions taken on an offer
are likewise surfaced through the platform's notification fabric, so the commercial pipeline is integrated with the rest of the
system rather than conducted out of band. The structured representation also means the state of a negotiation is always
explicit---which terms are on the table and where each offer stands---instead of being reconstructed from a thread of
correspondence.

\begin{figure}[htbp]
  \centering
  \includegraphics[width=0.95\textwidth]{fig-sponsors.png}
  \caption{The Sponsors module, presenting sponsor offers and contract negotiation as
  structured, auditable records.}
  \label{fig:sponsors}
\end{figure}

%======================================================================
\section{Reports and Analytics}
\label{sec:res-reports}

The Reports module, served by the reports-analytics service (port~8083), aggregates data from across the platform into
match analyses, team- and player-level analytics and training analytics (Figure~\ref{fig:reports}).

\paragraph{Cross-domain insight.} By unifying the previously siloed domains behind a single analytics layer, the module makes
cross-domain insight practical---for example relating match performance to squad availability or training load, a question
that is impossible to answer when each domain lives in its own spreadsheet. The service reads from the operational services'
data and presents the results through role-aware report views, so a performance analyst, a sport manager and a head coach each
see the analyses pertinent to their responsibilities. Because every domain shares the canonical identity scheme, a player's
match statistics, medical history and training attendance can be correlated into a single analytical picture without brittle
name-matching across systems.

\begin{figure}[htbp]
  \centering
  \includegraphics[width=0.95\textwidth]{fig-reports.png}
  \caption{The Reports and Analytics module, aggregating match, team, player and training
  data served by the reports-analytics service.}
  \label{fig:reports}
\end{figure}

%======================================================================
\section{Settings}
\label{sec:res-settings}

The Settings module gives every user control over their profile, preferences, notification choices and security, with changes
genuinely persisted rather than held only in the browser (Figure~\ref{fig:settings}).

\paragraph{Persisted preferences.} Profile, preference and notification settings are saved per user so that the experience is
consistent across sessions and devices. A user's notification choices govern which alerts they receive, and their preferences
shape the interface, both surviving a logout and a change of machine because they are stored server-side rather than in
\texttt{localStorage}.

\paragraph{Keycloak-verified password change.} The security tab supports a password change that is genuinely authoritative
across the platform, because it is mediated by the identity provider rather than by any local store. The change is a two-step
operation handled by the user-management service. First, the user's \emph{current} password is re-verified: the service performs
an OAuth password-grant against the \texttt{mscms} realm using the supplied username and current password, and treats the change
as authorised only if the grant returns an access token. A wrong current password causes the grant to fail with an
invalid-grant error, which the service interprets as a failed verification and rejects the change---so a stolen, unlocked
session cannot silently change the credential without knowing the existing password. Second, only after that verification
succeeds is the new password applied through the Keycloak administrative API. Because the credential lives in the identity
provider, the change takes effect uniformly everywhere the platform authenticates, with no local password copy to fall out of
sync.

\begin{figure}[htbp]
  \centering
  \includegraphics[width=0.95\textwidth]{fig-settings.png}
  \caption{The Settings module, with profile, preferences, notification and security tabs.
  Password changes are verified and applied through Keycloak.}
  \label{fig:settings}
\end{figure}

%======================================================================
\section{Analysis and Evaluation}
\label{sec:res-analysis}

Having presented the implemented modules, this section steps back to evaluate the delivered system as a whole: how completely it
satisfies the functional and non-functional requirements set out in Chapter~2, which design decisions proved most effective,
what the system's performance and user experience are like in practice, and how the implementation was tested and validated.

\subsection{Functional Requirements Coverage}
The implemented system satisfies the functional requirements across every domain identified in Chapter~2.
Authentication, role-based access and the role-aware dashboard realise the identity and onboarding
requirements~\cite{keycloak,jwt,rbac}. The squad, competition and match modules realise the player and
match-operations requirements, and in several respects exceed the original specification: the live-match engine's
broadcast-style phase machine, with knockout-only extra time and penalty shootouts and phase-bounded event minutes,
and the faithful reproduction of the modern Champions League format with its rotation-generated league phase, go beyond
the basic match recording originally envisaged. The medical module realises the medical-and-fitness requirements with a
complete, status-driven injury lifecycle from report to fitness-cleared recovery. The messages, notifications, training,
scouting, sponsor and reports modules realise the communication, notification and analytics requirements, with the
head-to-head, world-cup and trophy-room surfaces adding fan-facing context beyond the operational core. Crucially, the
consistent enforcement of role-based access at the identity provider, at the gateway and in the front-end
middleware---the three-layer scheme described in Chapter~3---means a permission rule applies uniformly whether a user
navigates the interface or calls the API directly, so the access control is not merely cosmetic.

\subsection{Non-Functional Requirements Coverage}
The non-functional requirements were met chiefly through architectural choices rather than after-the-fact tuning.
\emph{Maintainability and modularity} are served by the decomposition into six independently deployable services behind a
gateway and a service registry, each owning its own database, so a change to the medical schema cannot break the match
service. \emph{Security} is centralised in Keycloak and enforced at three layers, with no service trusting anything but a
validated token, and with the most sensitive operation---a password change---re-verifying the current credential before
proceeding. \emph{Scalability and loose coupling} follow from the event-driven notification pipeline: producers emit facts
to Kafka and remain ignorant of consumers, so the system can grow new reactions to existing events without modifying the
producers. \emph{Reliability} is supported by idempotent ingestion (de-duplicated by external identifier), the self-cleaning
Match Hub, and graceful degradation on the dashboard where one slow service suppresses only its own card. \emph{Usability}
is addressed by the single-origin API client, the automatic recovery on token expiry, and the timestamp-driven live clock and
countdowns that keep the interface responsive without server round-trips.

\subsection{Design Decisions That Worked Well}
Several decisions proved especially effective in practice. Making the back-end the single source of truth for external data,
with scheduled idempotent ingestion and read-through standings, kept the front-end simple and the data trustworthy while
respecting the data provider's free-tier constraints and never exposing the API key to the browser. Capturing the lineup
atomically with the match guaranteed that every match carries a reconstructable starting XI, which is precisely what makes the
live engine, the substitution logic and the match-details page meaningful rather than synthetic. Deriving the medical journey's
current stage solely from the injury status produced a workflow that is immune to leftover records and easy to reason about, and
allowing an in-match injury event to create a real medical record and withdraw the player demonstrated genuine, end-to-end
cross-service integration spanning the match, player and medical services in a single action. Resolving the knockout-versus-points
decision from the originating fixture's round, rather than from a loose match type, kept the live engine consistent with the
competition format. Finally, the event-driven notification fabric, fed by Kafka and the live engine, delivered the role-aware
alerting the requirements demanded with no manual intervention.

\subsection{Performance and User-Experience Observations}
The microservice decomposition allowed each domain to be developed and reasoned about in isolation, and routing every browser
request through a single gateway origin kept the API client and the cross-origin configuration straightforward~\cite{microservices,restapi}.
The live clock and countdown are computed from pure timestamp helpers, so the common case of simply advancing the clock incurs no
server round-trip and the live view stays responsive, while every state-changing action is persisted immediately so that a reload
faithfully reconstructs the match---including the on-pitch set, the consumed substitutions and the running phase. The dashboard's
parallel, per-card aggregation keeps perceived latency close to that of the slowest single dependency rather than the sum of all
of them. The self-cleaning Match Hub and the status-driven medical and squad logic markedly reduce the manual reconciliation that
burdens traditional tools, since a single status change propagates everywhere it is relevant. Two classes of defect encountered
during development are worth recording because their fixes shaped the design: a timezone-conversion bug that made matches appear
hours early was resolved by storing and parsing kickoff times as naive local wall-clock values, and a clock that could inflate into
nonsensical added time (an observed ``45+132''') was resolved by hard-capping the displayed minute at the phase length and showing
only admin-set stoppage. The principal real-world limitation observed is external rather than architectural: because the modelled
system date sits after the imported seasons' conclusion, the freely available data contains no upcoming or in-progress fixtures, so
live ingestion of genuinely new matches will populate only once a new season is published. The platform nonetheless demonstrates the
full live and ingestion machinery against the schedulable, club-created matches, and the ingestion pipeline will auto-fill league
fixtures when a new season appears in the provider's data.

\subsection{Testing and Validation Approach}
Validation combined automated and manual techniques appropriate to a system of this size. At the back-end, the domain services
carry controller-level tests---for example, tests around the match-event, match-formation and match-lineup controllers in the
training-match service and the injury and message controllers in the medical and notification services---that exercise the REST
surface and the entity mappings. The most behaviourally complex logic was deliberately structured as pure functions to make it
unit-testable in isolation: the live-clock and phase-machine helpers (phase derivation from a minute, phase-bounded recordable-minute
ranges, the stoppage-capped clock label, and the kickoff validator) take only timestamps and configuration and return deterministic
results, so their many edge cases---a minute past regulation, a knockout tie level at full time, a kickoff in the past or under the
five-minute lead---can be asserted directly. The fixture generators are likewise deterministic: the double round-robin and the
Champions League rotation algorithm produce the same fixture graph for the same input, which makes properties such as ``every team
plays exactly eight distinct opponents'' and ``no pairing repeats'' verifiable. Beyond automated tests, the system was validated by
end-to-end exercise of the user journeys against the running, containerised stack: provisioning a user through Keycloak and signing
in, creating a competition and confirming the generated fixtures, scheduling a match with a mandatory lineup, driving a live match
through its phases including a knockout tie into penalties, triggering an in-match injury and following the player through the
medical journey to a cleared fitness test, exchanging real-time chat messages over the WebSocket, and confirming that the resulting
Kafka-driven and live-engine alerts appeared in the notification centre. A noted constraint on this process was that automated
screenshot-based visual verification could not attach to the developer's running instance, so visual changes were validated by
inspecting routes, bundles and live-API behaviour and by manual review of the rendered interface, while behavioural correctness was
established through the automated tests and the end-to-end journeys described above.

\subsection{Summary}
Taken together, the modules presented in this chapter show that the platform delivers on its central promise: to consolidate the
fragmented operational concerns of a multi-sport club into a single, role-aware, cloud-native system in which match operations,
medical care, communication, recruitment, commercial activity and analytics are connected rather than siloed. The combination of a
real, idempotent ingestion pipeline, a behaviourally faithful live-match engine, a status-driven clinical injury workflow, a
real-time WebSocket communication channel and an event-driven notification fabric distinguishes the delivered system from the simple
CRUD application from which it evolved, and substantiates the design objectives stated at the outset of this report. The evaluation
above confirms that the functional and non-functional requirements of Chapter~2 are met or exceeded, and identifies the small number
of limitations---chiefly the post-season external data and the constraints on automated visual verification---that bound the present
implementation and motivate the future work discussed in the next chapter.
